\pdfoutput=1
\documentclass[11pt]{article}
\usepackage[T1]{fontenc}
\usepackage{lmodern}
\usepackage{amsmath,amssymb,amsthm,mathtools,booktabs,array}
\usepackage[margin=1in]{geometry}
\usepackage[colorlinks=true,urlcolor=blue,citecolor=blue]{hyperref}
\newtheorem{theorem}{Theorem}
\newtheorem{lemma}[theorem]{Lemma}
\newtheorem{corollary}[theorem]{Corollary}
\newtheorem{proposition}[theorem]{Proposition}
\theoremstyle{definition}
\newtheorem{definition}[theorem]{Definition}
\theoremstyle{remark}
\newtheorem{remark}[theorem]{Remark}
\DeclareMathOperator{\tw}{tw}
\DeclareMathOperator{\pw}{pw}
\DeclareMathOperator{\fvs}{fvs}
\newcommand{\XNLP}{\mathsf{XNLP}}
\newcommand{\W}[1]{\mathsf{W}[#1]}
\newcommand{\Z}{\mathbb{Z}}
\newcommand{\Hom}{\operatorname{Hom}}
\newcommand{\lean}[1]{\nolinkurl{#1}}
\makeatletter\g@addto@macro\UrlBreaks{\do\_}\makeatother
\emergencystretch=2em
\title{Homomorphism to an odd cycle is XNLP-complete\newline
parameterized by pathwidth}
% Author block follows the companion manuscript in euclid-uniform-processes.
\author{Jason Davies\thanks{\href{mailto:jason@jasondavies.com}{\nolinkurl{jason@jasondavies.com}}}}
\date{Draft --- November 16, 2026}

\begin{document}
\maketitle

\begin{abstract}
We prove that deciding whether a finite simple undirected graph admits a
homomorphism to an odd cycle given in the input is XNLP-complete when a path
decomposition of the source is supplied and its width is the parameter.
The reduction from Chained Multicoloured Clique with $k$ colours produces
source pathwidth at most $2k+4$ and treewidth at most $2k+3$, independently of
the number of layers. Consequently, graph homomorphism parameterized by
the maximum of the two input treewidths is XNLP-hard and W[$t$]-hard for
every fixed positive $t$, even when the target has treewidth two.
A simpler reduction proves W[1]-hardness parameterized by the feedback
vertex number of the source. Both constructions use only copies of an odd
cycle and paths: cycle automorphisms encode choices, and paths of length
$n-2$ enforce distinct images in $C_n$.
Decision equivalence, the treewidth bounds, and polynomial size bounds for
the two reductions are checked in Lean~4. The pathwidth bounds and the
complexity-theoretic conclusions are proved in the text.
\end{abstract}

\noindent\textbf{AI transparency statement.}
The mathematical development and Lean formalisation were carried out with
AI assistance under human direction. This manuscript draft was prepared
by an AI assistant. The scope of formal verification is stated in
Appendix~\ref{app:lean}; further disclosure appears at the end of the paper.

\section{Introduction}\label{sec:intro}

A graph homomorphism from $G$ to $H$ is a map on vertices that preserves
adjacency. When a tree decomposition of $G$ of width $w$ is available,
the usual dynamic programme decides whether such a map exists in time
$|V(H)|^{O(w+1)}$ times a polynomial in the input size. This is polynomial
for every fixed $w$, but the exponent depends on $w$.
Amarilli asked whether bounding the treewidth of the target as well makes
the problem fixed-parameter tractable: is homomorphism parameterized by
$\max\{\tw(G),\tw(H)\}$ FPT or W[1]-hard~\cite{amarilli2016}?

We give an explicit hardness construction whose target is a single odd
cycle. Its length varies with the input; it is not a fixed target and is
not included in the parameter. Thus the target has treewidth and pathwidth
exactly two, while the source width is bounded by a function of the
parameter of the starting problem.

\begin{theorem}\label{thm:main}
Given a finite simple undirected graph $G$, an odd integer $n\geq5$ in unary,
and a path decomposition of $G$ of width at most $w$, deciding whether
$G\to C_n$ is XNLP-complete parameterized by $w$.
\end{theorem}

We use parameterized logspace reductions. As in the standard treatment of
pathwidth-parameterized XNLP problems~\cite[Section~2.4]{bgns2024}, the
decomposition is supplied. Membership without this extra input is not
established here.

\begin{theorem}\label{thm:plain}
Homomorphism to an input odd cycle is W[1]-hard parameterized by the
feedback vertex number of the source. A reduction from Multicoloured
Clique with $k$ classes produces a graph $G$ satisfying
\[
 \fvs(G)\leq k+1,\qquad \tw(G)\leq k+2,\qquad \pw(G)\leq k+3.
\]
\end{theorem}

\begin{corollary}\label{cor:both}
Graph homomorphism parameterized by $\max\{\tw(G),\tw(H)\}$ is XNLP-hard
and W[$t$]-hard for every fixed positive integer $t$, even when $H$ is an
odd cycle. The same hardness holds with pathwidth in place of treewidth.
\end{corollary}

The target bound two is best possible for hardness of this form unless
FPT${}={}$W[1]. For a nonempty forest $H$ containing an edge, $G\to H$ holds
exactly when $G$ is bipartite: map its two parts to the endpoints of any
edge of $H$. If $H$ is nonempty and edgeless, the condition is that $G$ be
edgeless; if $H$ is empty, $G$ must be empty. These cases are polynomial-time
decidable.

\paragraph{Construction.}
Every endomorphism of an odd cycle $C_n$ is a rotation or reflection.
We use a source copy of $C_n$ as a \emph{dial}: the image of its designated
vertex records a choice. A path of length $n-2$ between two source vertices
can be mapped into $C_n$ exactly when their images are different.
Adding such paths forbids individual values and differences of values.
An elementary labelling makes a permitted difference determine exactly
one pair of choices, even in the presence of reflections.

For a chain of instances, a single global reference cycle would obstruct
the intended width argument. We instead give each choice its own reference
cycle, with all reference cycles sharing the same edge. Their maps to the
target necessarily agree in coordinates. This lets us process the layers
in order, keeping only two layers of choice hubs and the shared edge.

\paragraph{Prior work and attribution.}
The W[1]-hardness answer to Amarilli's question was already announced in
his comment of May~30, 2025~\cite{amarilli2016}. We make no priority claim
for that answer. The relationship between the construction here and those
unpublished reductions has not been established.
Our XNLP-hardness starting point is the Chained Multicoloured Clique
theorem of Bodlaender, Groenland, Nederlof and
Swennenhuis~\cite[Theorem~11]{bgns2024}.
Related fine-grained work by Piecyk and Rz\k{a}\.{z}ewski studies list
homomorphisms parameterized by feedback vertex number and cutwidth, and
also obtains lower bounds for ordinary homomorphisms to odd cycles
parameterized by cutwidth~\cite{pr2021}. Those bounds concern the base of
an exponential running time; the present statements concern uniform
parameterized hardness with the cycle included in the input.

\section{Preliminaries}\label{sec:prelim}

All graphs are finite, simple, undirected and explicitly represented.
We write $G\to H$ when there is a homomorphism from $G$ to $H$.
For odd $n\geq3$, the graph $C_n$ has vertex set $\Z/n\Z$ and edges
$\{x,x+1\}$. Arithmetic in coordinates is modulo $n$ unless stated
otherwise. A walk may repeat vertices and edges.

A tree decomposition consists of bags indexed by a tree, covering every
vertex and every edge, such that the bags containing a given vertex induce
a connected subtree. Its width is the largest bag size minus one.
A path decomposition uses a path as its indexing tree; equivalently, the
bags containing any fixed vertex form an interval in the bag sequence.
The minima of these widths are $\tw(G)$ and $\pw(G)$. The feedback vertex
number $\fvs(G)$ is the minimum size of a vertex set whose deletion leaves
a forest.

The class $\XNLP$ consists of parameterized problems decidable in
nondeterministic time $f(k)L^{O(1)}$ and space $f(k)\log L$, simultaneously,
where $L$ is the input length and $f$ is computable. Parameterized logspace
reductions use $g(k)+O(\log L)$ work space and bound the output parameter by
a computable function of $k$. Our reductions in fact use ordinary
logarithmic space and polynomial time. XNLP-hardness implies W[$t$]-hardness
for each fixed positive $t$~\cite[Lemma~2]{bgns2024}.

In \emph{Multicoloured Clique}, the vertices of an input graph are
partitioned into $k$ classes and one asks for a clique choosing one vertex
per class. Its W[1]-hardness follows directly from that of Clique: replace
each vertex by one copy in each of $k$ classes, and join copies in distinct
classes exactly when the original vertices are adjacent. An empty class
is a no-instance. The case $k=1$ is directly decidable. In the remaining
cases, padding classes by isolated vertices gives a common class size
$N\geq1$ without changing the answer.

In \emph{Chained Multicoloured Clique}, there are $r$ layers, each
partitioned into $k$ colour classes. A solution chooses one vertex per
class so that the chosen vertices of every consecutive pair of layers
form a clique. This problem is XNLP-complete parameterized by
$k$~\cite[Theorem~11]{bgns2024}. We use $r\geq2$, $k\geq1$; the omitted
degenerate cases are directly decidable. Empty classes give no-instances,
and isolated vertices again equalise all class sizes to $N\geq1$.

\section{Two odd-cycle gadgets}\label{sec:gadgets}

\begin{lemma}[Dials]\label{lem:dial}
Every homomorphism $C_n\to C_n$, for odd $n\geq3$, has the form
\[
 x\longmapsto r+\varepsilon x,
 \qquad r\in\Z/n\Z,\quad \varepsilon\in\{1,-1\}.
\]
\end{lemma}
\begin{proof}
The image of the source cycle is a closed walk of length $n$ in the target.
Each step has displacement $+1$ or $-1$. If $s$ steps are negative, the
total displacement is $n-2s$, a multiple of $n$ in $[-n,n]$. It is odd,
so it is $n$ or $-n$. Thus every step has the same sign.
\end{proof}

\begin{lemma}[Links]\label{lem:link}
For odd $n\geq3$, two vertices $x,y$ of $C_n$ are joined by a walk of
length $n-2$ if and only if $x\ne y$.
\end{lemma}
\begin{proof}
The possible displacements are
\[
 -(n-2),\;-(n-4),\;\ldots,\;-1,\;1,\;\ldots,\;n-4,\;n-2.
\]
Modulo $n$, the positive terms give all odd nonzero residues and the
negative terms give all even nonzero residues. These are exactly the
$n-1$ nonzero residues. Each displacement is realised by taking the
appropriate number of positive and negative steps.
\end{proof}

We call a path of length $n-2$ with fresh internal vertices a \emph{link}.
By Lemma~\ref{lem:link}, attaching a link between two vertices imposes
exactly the constraint that their images be distinct. All links in what
follows have pairwise disjoint interiors, also disjoint from the dials.

\section{Labels that distinguish ordered pairs}\label{sec:labels}

Let $K,N\geq1$, and index $K$ classes by $s\in\{0,\ldots,K-1\}$.
For the $a$th vertex of class $s$, with $1\leq a\leq N$, set
\begin{equation}\label{eq:labels}
 Q=2NK+1,\qquad p_s=Q+s,\qquad \lambda_s(a)=a p_s,
 \qquad M=N(Q+K)+2,\qquad n=2M+1.
\end{equation}
In particular $n\geq5$ is odd, $1\leq\lambda_s(a)\leq M-2$, and
$n=O(N^2K)$.

\begin{lemma}[Signed differences]\label{lem:labels}
Fix $s<u$ and $a,a',b,b'\in\{1,\ldots,N\}$. If
\[
 \delta\bigl(\lambda_u(b')-\lambda_s(a')\bigr)
 \equiv \lambda_u(b)-\lambda_s(a)\pmod n,
 \qquad \delta\in\{1,-1\},
\]
then $\delta=1$, $a'=a$ and $b'=b$.
\end{lemma}
\begin{proof}
Each side has absolute value less than $M$, so the congruence is an
equality of integers. We first observe that
\begin{equation}\label{eq:arithmetic}
 x p_u=y p_s\quad\text{and}\quad |ys-xu|<Q
 \quad\Longrightarrow\quad x=y=0.
\end{equation}
Indeed, $(x-y)Q=ys-xu$, whence $x=y$; then $x(u-s)=0$.

For $\delta=1$, take $x=b'-b$ and $y=a'-a$. Both have absolute value
at most $N-1$, and
$|ys-xu|\leq2(N-1)(K-1)<Q$.
Equation~\eqref{eq:arithmetic} gives $a'=a$ and $b'=b$.
For $\delta=-1$, take $x=b'+b$ and $y=a'+a$. Now $2\leq x,y\leq2N$,
so both $ys$ and $xu$ lie in $[0,2N(K-1)]$; again $|ys-xu|<Q$.
Equation~\eqref{eq:arithmetic} would give $x=0$, a contradiction.
\end{proof}

For an input graph on these classes, write
\begin{align}
 \Lambda_s&=\{\lambda_s(a):1\leq a\leq N\},\label{eq:lambda}\\
 D_{su}&=\{\lambda_u(b)-\lambda_s(a)\bmod n:
                \text{the vertices $(s,a)$ and $(u,b)$ are adjacent}\}.
 \label{eq:differences}
\end{align}
Lemma~\ref{lem:labels} shows that membership in $D_{su}$ identifies an
actual input edge, and cannot be satisfied by reversing a different pair
of choices.

\section{The plain reduction}\label{sec:plain}

Apply~\eqref{eq:labels} with $K=k$ to a Multicoloured Clique instance.
Take a reference dial $R=C_n$ and disjoint choice dials
$F_0,\ldots,F_{k-1}$, also copies of $C_n$. Let $r_d$ denote coordinate
$d$ of $R$, and $f^s_d$ coordinate $d$ of $F_s$. Its hub is $h_s=f^s_0$.
Construct $G$ by adding the following links:
\begin{enumerate}
\item for each $s$ and each $d\notin\Lambda_s$, a link from $r_d$ to $h_s$;
\item for each $s<u$ and each $d\notin D_{su}$, a link from $f^s_d$ to $h_u$.
\end{enumerate}
The target is $C_n$.

\begin{proposition}\label{prop:plain-correct}
The constructed graph $G$ admits a homomorphism to $C_n$ if and only if
the input has a multicoloured clique.
\end{proposition}
\begin{proof}
Given a clique selecting vertex $a_s$ in class $s$, map $r_d$ to $d$ and
$f^s_d$ to $\lambda_s(a_s)+d$. These maps preserve the dial edges.
A class link has distinct endpoint images since $d\notin\Lambda_s$.
An edge link also has distinct endpoint images: equality would imply
$d=\lambda_u(a_u)-\lambda_s(a_s)\in D_{su}$. Every link therefore extends
by Lemma~\ref{lem:link}.

Conversely, suppose $\phi:G\to C_n$. By Lemma~\ref{lem:dial} and a target
automorphism, we may assume $\phi(r_d)=d$ for every $d$. The class links
exclude every residue outside $\Lambda_s$, so
$\phi(h_s)=\lambda_s(a_s)$ for some $a_s$. On $F_s$ the map has the form
\[
 \phi(f^s_d)=\lambda_s(a_s)+\varepsilon_s d,
 \qquad \varepsilon_s\in\{1,-1\}.
\]
For $s<u$, all edge links can map only if
\begin{equation}\label{eq:allowed-difference}
 \varepsilon_s\bigl(\lambda_u(a_u)-\lambda_s(a_s)\bigr)\in D_{su}.
\end{equation}
To see this, the residue on the left is the unique $d$ for which
$\phi(f^s_d)=\phi(h_u)$; it cannot be a forbidden residue.
By~\eqref{eq:differences} and Lemma~\ref{lem:labels}, the chosen vertices
in classes $s$ and $u$ are adjacent. This holds for every pair of classes.
\end{proof}

\begin{proposition}\label{prop:plain-width}
The graph $G$ has feedback vertex number at most $k+1$, treewidth at most
$k+2$, and pathwidth at most $k+3$.
\end{proposition}
\begin{proof}
Delete the coordinate-zero vertex of each of the $k+1$ dials. Each dial
becomes a path. Every link loses its hub endpoint, leaving a pendant path
at its other endpoint, or a separate path if both endpoints were deleted.
The remaining graph is a forest, proving the feedback vertex bound.
Add the deleted vertices to every bag of a width-one decomposition of
this forest to obtain the treewidth bound.

Each component of the forest is a path with pendant paths attached, or
is itself a path. For a path decomposition, traverse the main path and
process the pendant paths at each vertex one at a time. Retain the current
main-path vertex and at most two consecutive vertices of the pendant path.
Between successive main-path vertices use a two-vertex bag. This has width
at most two. Concatenate the component decompositions and add the $k+1$
deleted vertices to all bags, obtaining width at most $k+3$.
\end{proof}

The number of vertices is at most
\[
 (k+1)n+\left(k+\binom{k}{2}\right)n(n-3)\leq(k+1)^2n^2.
\]
All links and their interiors can be enumerated in polynomial time.
Propositions~\ref{prop:plain-correct} and~\ref{prop:plain-width} therefore
prove Theorem~\ref{thm:plain}.

\section{The chained reduction}\label{sec:chained}

Consider a Chained Multicoloured Clique instance with $r\geq2$ layers and
$k\geq1$ colours per layer. Index its $K=rk$ classes, or \emph{slots},
by $s=0,\ldots,K-1$, with layer $\lfloor s/k\rfloor$.
A pair $s<u$ is \emph{required} when
\begin{equation}\label{eq:required}
 \left\lfloor\frac{u}{k}\right\rfloor
 \leq \left\lfloor\frac{s}{k}\right\rfloor+1.
\end{equation}
Use~\eqref{eq:labels} with these $K$ slots.

\subsection{Synchronised reference cycles}

For each slot $s$, take a reference dial $R_s$ and a choice dial $F_s$.
Identify coordinate zero of every $R_s$ to one vertex $A$, and coordinate
one of every $R_s$ to one vertex $B$. Thus the reference cycles share
exactly the edge $AB$. Their remaining vertices, and all choice cycles,
are disjoint. Write $r^s_d$ and $f^s_d$ for their coordinates and
$h_s=f^s_0$ for a choice hub.

Add a class link from $r^s_d$ to $h_s$ for every $d\notin\Lambda_s$.
For each required pair $s<u$, add an edge link from $f^s_d$ to $h_u$ for
every $d\notin D_{su}$. There are no edge links for nonrequired pairs.
All interiors are fresh. Denote the resulting graph by $G$.

\begin{proposition}\label{prop:chained-correct}
The graph $G$ admits a homomorphism to $C_n$ if and only if the input has
a chained multicoloured clique.
\end{proposition}
\begin{proof}
In any homomorphism, the adjacent vertices $A,B$ have adjacent images.
Compose with a target automorphism so that their images are $0,1$.
The restriction to each reference cycle is an automorphism by
Lemma~\ref{lem:dial}; these two images force that automorphism to be the
coordinate identity. Hence $\phi(r^s_d)=d$ simultaneously for all $s,d$.

Class links force each choice hub to a vertex label. For every required
pair, the argument establishing~\eqref{eq:allowed-difference} applies
unchanged. Lemma~\ref{lem:labels} forces the selected pair to be an input
edge. Thus all required pairs are edges and the selection is a chained
multicoloured clique.

Conversely, map each reference cycle by identity and each choice cycle by
rotation through its selected label. The selected vertices are adjacent
for every required pair, so the same endpoint argument as in
Proposition~\ref{prop:plain-correct} extends the map over all links.
\end{proof}

\subsection{Width bounds}\label{sec:chained-width}

For layer $t\in\{0,\ldots,r-1\}$, retain the set
\begin{equation}\label{eq:retained}
 S_t=\{A,B\}\cup
 \{h_s:\lfloor s/k\rfloor\in\{t,t+1\}\}.
\end{equation}
Then $|S_t|\leq2k+2$. All vertices other than $A,B$ and the choice hubs
are assigned to one phase: the private vertices of $R_s,F_s$ belong to
phase $\lfloor s/k\rfloor$, as do the internal vertices of every link
starting on one of those cycles.

\begin{proposition}\label{prop:chained-width}
The chained source satisfies $\tw(G)\leq2k+3$ and $\pw(G)\leq2k+4$.
Both decompositions can be explicitly constructed.
\end{proposition}
\begin{proof}
After removing all retained vertices, the remaining graph is a forest.
Each reference cycle has lost $A,B$, each choice cycle has lost its hub,
and each link has lost its destination hub. Its components lie within
single phases. Moreover, a nonretained vertex in phase $t$ has retained
neighbours only in $S_t$, by~\eqref{eq:required}.

For a tree decomposition, make a path of bags $S_0,\ldots,S_{r-1}$.
For each component of the nonretained forest, root it arbitrarily and
attach its usual width-one decomposition to its phase bag, adding $S_t$
to every bag of that component. The resulting bags have size at most
$2k+4$. Every nonretained edge is covered within its component, every
mixed edge is covered by the added $S_t$, and the only possible edge
between retained vertices is $AB$. Retained vertices occur in intervals
of the phase path and in the trees attached to those intervals; all
occurrence sets are connected. This proves the treewidth bound.

For a path decomposition, process phases in increasing order, retaining
$S_t$ throughout phase $t$. Traverse each reference and choice cycle of
that layer in coordinate order. At its current vertex $x$, process all
links designated to start at that occurrence of $x$, one at a time.
Their destination hubs belong to $S_t$. For a link with vertices
$x,z_1,\ldots,z_{n-3},h$, use bags
\[
 S_t\cup\{x,z_1\},\qquad
 S_t\cup\{x,z_i,z_{i+1}\},\qquad
 S_t\cup\{x,z_{n-3},h\}.
\]
Use $S_t\cup\{x,x_{\rm next}\}$ for consecutive cycle edges. The closing
edge of a choice cycle meets its retained hub; both reference anchors
are retained. At a phase boundary, forget the older hubs before
introducing the following layer's hubs; the common hubs remain present.

Every edge is covered. A private link vertex appears during its traversal,
and a private cycle vertex appears during its visit and the adjacent
cycle-edge bags. Each hub occurs in one or two consecutive phases, while
$A,B$ occur throughout. Thus vertex occurrences are intervals.
The largest bags have at most $|S_t|+3\leq2k+5$ vertices, proving the
pathwidth bound.
\end{proof}

\subsection{Reduction resources}\label{sec:resources}

There are $K$ reference cycles and $K$ choice cycles. Their union has
$2+K(n-2)+Kn$ vertices. The number of required pairs is
\[
 r\binom{k}{2}+(r-1)k^2.
\]
Together with the $K=rk$ class constraints, this gives $O(rk^2n)$ links
and $O(rk^2n^2)$ total vertices and edges. Since $n=O(N^2rk)$, this size
is polynomial in the input size. A convenient coarser vertex bound,
used in the formalisation, is $2(K+1)^2n^2$.

The graph and its path decomposition can be emitted in logarithmic space.
Use a constant number of counters for slots, residues and link positions.
All these values have polynomial magnitude in the explicit input size.
Labels are computed by integer arithmetic. To test $d\in D_{su}$, rescan
the input edges and test their label differences; no set of differences
is stored. Identify output vertices by tuples of the generating indices,
or recompute sequential numbers by rescanning the same enumeration.
Output each bag of Proposition~\ref{prop:chained-width} vertex by vertex;
in particular, enumerate its retained hubs instead of storing the bag.
Only a constant number of nested polynomial loops and input scans is
required, so the time is polynomial as well.

An empty input class is mapped to the fixed no-instance $K_3\to C_5$,
with its width-two source decomposition. This is a no-instance because
$C_5$ has no triangle. Padding and indexing nonempty classes can be
performed by input scans in the same space bound. The output parameter is
at most $2k+4$, independently of $r$ and $N$.

Propositions~\ref{prop:chained-correct} and~\ref{prop:chained-width} now give
XNLP-hardness in Theorem~\ref{thm:main}. They also give the hardness
claims of Corollary~\ref{cor:both}, using
\cite[Lemma~2]{bgns2024}. A target decomposition of width two is available
explicitly, with bags $\{0,d,d+1\}$ for $d=1,\ldots,n-2$.

\section{Membership in XNLP}\label{sec:membership}

We give the membership argument for an arbitrary explicitly represented
target graph $H$; it applies in particular to $C_n$ with $n$ given in unary.
Let $L\geq2$ be the total input length, and let $B_1,\ldots,B_\ell$ be a
supplied path decomposition of width at most $w$, with $\ell\geq1$.
Use vertex identifiers that index the explicit vertex lists.

\begin{lemma}\label{lem:gluing}
There is a homomorphism $G\to H$ if and only if there are assignments
$a_i:B_i\to V(H)$ such that each $a_i$ preserves source edges within its
bag, and $a_i,a_{i+1}$ agree on $B_i\cap B_{i+1}$ for every $i$.
\end{lemma}
\begin{proof}
A homomorphism restricts to such assignments. Conversely, a vertex occurs
in an interval of bags, so agreement across consecutive overlaps makes
all its assigned images equal. Vertex coverage therefore defines one
map on $V(G)$. Each edge is contained in a bag, where the local edge
condition proves that its endpoints map to adjacent vertices.
\end{proof}

Process the bags from left to right. Guess an image for each vertex of the
current bag, retaining the preceding bag's assignments. Check agreement
on the overlap and check all edges whose endpoints lie in the current bag.
Then discard the preceding assignments. By Lemma~\ref{lem:gluing}, an
accepting branch exists exactly when a homomorphism exists.
The algorithm never constructs a table of all assignments.

There are at most $2(w+1)$ stored vertex-image pairs, each using
$O(\log L)$ bits. Counters and adjacency checks require an additional
$O(\log L)$ bits. The work space is consequently
$O((w+1)\log L)$. Guess an image using a binary index and reject indices
outside the target vertex range. As all bags and edges are explicitly
listed, repeated scans and comparisons take polynomial time with an
exponent independent of $w$.

This may be formulated as a total decision problem rather than a promise
problem. Reject malformed inputs and invalid decompositions first.
For each vertex, scan the bag sequence and check that its occurrence set
is a nonempty interval. For each edge, find a bag containing both
endpoints. Check valid vertex identifiers, no repeated entries in a bag,
and the width bound. These checks need only logarithmic space and
polynomial time: a constant number of counters and an occurrence flag
suffice for each repeated scan. An oversized binary width bound can be
capped at $L$, since no explicitly listed bag has more than $L$ entries.
For an empty source graph, allow one empty bag.

Thus the time and space bounds required for XNLP hold simultaneously.
This completes the proof of Theorem~\ref{thm:main}.

\section{Scope and further questions}\label{sec:scope}

The main theorem assumes a supplied path decomposition. The reduction
also proves XNLP-hardness when that decomposition is omitted, but the
membership argument does not find one within the required resources.
Similarly, the treewidth result is a hardness statement; we do not claim
XNLP-completeness parameterized by treewidth.

The plain reduction has a feedback vertex set of size $k+1$. The chained
construction does not preserve that stronger restriction: its $rk$
choice cycles are vertex-disjoint, so any feedback vertex set has size
at least $rk$. Accordingly Theorem~\ref{thm:plain} remains a separate
structural refinement of the main result.

For each fixed source width, the usual dynamic programme remains
polynomial. In particular, these results do not assert hardness when
both source and target have pathwidth two. Only the target width is fixed
at two; the source bound grows with the source parameter.

\appendix
\section{Formal verification}\label{app:lean}

The development uses Lean~4.35.0~\cite{deMouraUllrich2021} and the pinned
mathlib release v4.35.0~\cite{mathlib2020}. The targets in the final
reduction statements are mathlib's \lean{SimpleGraph.cycleGraph}.
The repository defines a tree-decomposition certificate by vertex
coverage, edge coverage and connected occurrence bags; \lean{HasWidth}
asserts the existence of such a certificate with bounded bag size.
Table~\ref{tab:lean} records the main correspondence.

\begin{table}[ht]
\centering\small
\caption{Principal checked statements. Names are relative to
\texttt{TreewidthHom}.}
\label{tab:lean}
\begin{tabular}{@{}>{\raggedright\arraybackslash}p{0.34\textwidth}>{\raggedright\arraybackslash}p{0.62\textwidth}@{}}
\toprule
Mathematical statement & Lean declaration(s)\\\midrule
Dial rigidity, Lemma~\ref{lem:dial} & \lean{OddCycle.dial}\\
Link constraint, Lemma~\ref{lem:link} & \lean{OddCycle.walk_iff_ne}, \lean{linked_hom_iff}\\
Signed-difference decoding & \lean{OddCycle.label_solve}\\
Plain correctness, treewidth and size & \lean{OddCycle.reduction_cycleGraph}\\
Chained correctness, treewidth and size & \lean{OddCycle.Chained.reduction}\\
Local-to-global membership lemma & \lean{Decomposition.hom_iff_pathBagAssignments}\\
Two-bag record bound & \lean{Decomposition.two_bag_records_le}\\
\bottomrule
\end{tabular}
\end{table}

The plain theorem gives source treewidth at most $k+2$, target treewidth
at most two, and at most $(k+1)^2n^2$ source vertices. The chained theorem
gives source treewidth at most $2k+3$ and at most $2(K+1)^2n^2$ source
vertices, for any number $K$ of slots and any positive layer size $k$.
Taking $K=rk$ gives the instances used here. The label-solving theorem
states equality of the decoded choices; its proof eliminates the negative
orientation as in Lemma~\ref{lem:labels}.

The pathwidth bounds are proved in this paper but are not stated by the
Lean reduction theorems. Nor does Lean state the feedback vertex number
bound as a separate theorem. It checks the parent-pointer forest
construction used for the corresponding treewidth proof.
The membership lemmas formalise the consistency argument and the number
of stored vertex-image records; they do not implement a machine model.
The published hardness starting points, polynomial-time and logspace
resource analyses, and parameterized-complexity conclusions remain
external to Lean.

The axiom audit allows only \lean{propext}, \lean{Classical.choice} and
\lean{Quot.sound}, and rejects proof admissions. The development can be
built and audited from the repository root by
\begin{verbatim}
lake build
python3 scripts/check_formalization.py
\end{verbatim}
The precise theorem map and trust boundary are recorded in
\lean{FORMALIZATION.md}; machine-readable audit output is in
\lean{sources/lean-verification.json}.

\section{Finite verification}\label{app:finite}

The repository also implements both reductions in Python. For each test,
it compares direct clique search with a graph-homomorphism search that
uses arc consistency and backtracking on dial vertices. After these
vertices are assigned, the remaining graph is a forest. Every returned
homomorphism is checked on all edges.

The recorded plain suite covers 4,382 instances, including 1,798
no-instances. The recorded chained suite covers 361 instances, including
124 no-instances, with up to 1,265,582 source vertices. On each chained
instance it checks an explicit path decomposition for vertex coverage,
edge coverage, interval occurrences and width at most $2k+4$. One test
has feasible consecutive layer pairs individually but no choice in the
middle layer that works for both; it is correctly rejected.
These finite checks supplement the general proofs and do not establish
their correctness on arbitrary inputs.

\begin{verbatim}
python3 scripts/run_odd_cycle_checks.py
python3 scripts/run_odd_cycle_chained_checks.py
\end{verbatim}
The corresponding reports are:
\begin{verbatim}
sources/odd-cycle-results.json
sources/odd-cycle-chained-results.json
\end{verbatim}

\section*{Declaration of AI-assisted research and writing}

AI systems assisted in developing the reductions, writing and checking
the Lean formalisation, and preparing this manuscript under the direction
of the human author. This initial manuscript draft was written by an AI
assistant. This disclosure does not assert that the draft has received
independent human review. Kernel checking supplies the guarantees stated
in Appendix~\ref{app:lean}; it does not certify novelty or the mathematical
claims left outside the formal development.

\begin{thebibliography}{10}
\raggedright

\bibitem{amarilli2016}
Antoine Amarilli.
Complexity of the homomorphism problem parameterized by treewidth.
\emph{Theoretical Computer Science Stack Exchange}, question~34877,
June~2, 2016; comment announcing W[1]-hardness, May~30, 2025.
\url{https://cstheory.stackexchange.com/q/34877}.

\bibitem{bgns2024}
Hans L. Bodlaender, Carla Groenland, Jesper Nederlof and C\'eline Swennenhuis.
Parameterized problems complete for nondeterministic FPT time and
logarithmic space.
\emph{Information and Computation} 300:105195, 2024.
\href{https://doi.org/10.1016/j.ic.2024.105195}{doi:10.1016/j.ic.2024.105195}.

\bibitem{deMouraUllrich2021}
Leonardo de Moura and Sebastian Ullrich.
The Lean~4 theorem prover and programming language.
In \emph{Automated Deduction -- CADE~28}, Lecture Notes in Computer Science,
vol.~12699, pp.~625--635. Springer, 2021.
\href{https://doi.org/10.1007/978-3-030-79876-5_37}{doi:10.1007/978-3-030-79876-5\_37}.

\bibitem{mathlib2020}
The mathlib Community.
The Lean mathematical library.
In \emph{Proceedings of the 9th ACM SIGPLAN International Conference on
Certified Programs and Proofs (CPP 2020)}, pp.~367--381. ACM, 2020.
\href{https://doi.org/10.1145/3372885.3373824}{doi:10.1145/3372885.3373824}.

\bibitem{pr2021}
Marta Piecyk and Pawe{\l} Rz\k{a}\.{z}ewski.
Fine-grained complexity of the list homomorphism problem: feedback vertex
set and cutwidth.
In \emph{38th International Symposium on Theoretical Aspects of Computer
Science (STACS 2021)}, LIPIcs, vol.~187, pp.~56:1--56:17.
Schloss Dagstuhl, 2021.
\href{https://doi.org/10.4230/LIPIcs.STACS.2021.56}{doi:10.4230/LIPIcs.STACS.2021.56}.

\end{thebibliography}
\end{document}
